\section{МЕТОДЫ}

Конструирование и взвешивание сетей проводили при помощи пакета genelens v.1.0.10. Для выгрузки генов сигнальных путей использовали \texttt{participants\_reference\_entities} пакета reactome2py v.3.0.0. Выгрузка генов-мишеней микроРНК из miRTarBase реализована в модуле \texttt{Targets} пакета genelens.

\subsection{Представление данных}
Каждый сигнальный путь представляется в виде неориентированного графа $G = (V_G, E_G)$ (сети), где:
\begin{itemize}
    \item $V_G$ --- множество вершин, соответствующих генам сигнального пути;
    \item $E_G \subseteq V_G \times V_G$ --- множество рёбер, соответствующих взаимодействиям белковых продуктов.
\end{itemize}

Аналогично, микроРНК представляется как граф $M = (V_M, E_M)$ её генов-мишеней. Оба графа извлекаются из человеческого интерактома как его подграфы.

\subsection{Взвешивание вершин}
Каждой вершине $v \in V_G \cup V_M$ сопоставляется нормированный вес $w(v) \in [0,1]$, вычисляемый как сумма betweenness centrality (BC) и degree centrality:

\begin{equation}
w(v) = \frac{\operatorname{BC}(v) + \deg(v)}{\max_{u \in V_G \cup V_M} (\operatorname{BC}(u) + \deg(u))}
\end{equation}

\subsection{Вычисление расстояния между графами}
Расстояние $D(G,M)$ между графами сигнального пути $G$ и микроРНК $M$ вычисляется по формуле:

\begin{equation}
D(G,M) = \sum_{v \in V_G \cap V_M} |w_G(v) - w_M(v)| + \alpha \sum_{v \in V_G \setminus V_M} w_G(v) + \beta \sum_{v \in V_M \setminus V_G} w_M(v)
\label{eq:optim_task}
\end{equation}

где:
\begin{itemize}
    \item $\alpha = 1$ --- штраф за гены сигнального пути, не покрытые микроРНК
    \item $\beta = \frac{1}{2}$ --- штраф за избыточные гены-мишени микроРНК
\end{itemize}

\subsection{Отбор значимых микроРНК}
Для выбора оптимального порога $\tau$ использовали анализ функции распределения расстояний:

\begin{enumerate}
    \item Строили эмпирическую функцию распределения (ECDF):
    \begin{equation}
    F(x) = \frac{|\{M_i | D(G,M_i) \leq x\}|}{N}
    \end{equation}
    где $N$ - общее количество микроРНК
    
    \item Вычисляли локальный рост функции:
    \begin{equation}
    \Delta F(x) = F(x) - F(x - \Delta x)
    \end{equation}
    с шагом $\Delta x = 0.05$
    
    \item Определяли порог $\tau$ как минимальное $x$, при котором:
    \begin{equation}
    \Delta F(x) > 0.2
    \end{equation}
    что соответствует резкому увеличению числа микроРНК в окрестности точки
\end{enumerate}

\subsection{Оптимизационный отбор}
Среди микроРНК с $D(G,M) \leq \tau$ решали задачу лексикографической оптимизации:
\begin{equation*}
    \begin{cases}
    1) \max |V_G \cap V_{M_i}| \\
    2) \min D(G,M_i) \text{ при равном покрытии}
    \end{cases}
\end{equation*}

что эквивалентно выбору микроРНК с:
\begin{itemize}
    \item Максимальным покрытием генов сигнального пути
    \item Минимальным расстоянием до порога $\tau$
\end{itemize}

Для анализа использовали только гены, экспрессирующиеся в ткани сердца по данным Human Protein Atlas \cite{osmak2020mirna}. Все вычисления реализованы в пакете GeneLens. Код доступен на GitHub:

\url{https://github.com/GJOsmak/miRNA_Assignment_problem}